\section{Euclid's algorithm: meaning and termination}
A positive integer $d$ divides an integer $a$ if $a=dk$ for some integer $k$. Write $d\mid a$. For nonnegative integers $a,b$ not both zero, their greatest common divisor, $\gcd(a,b)$, is the largest positive integer dividing both. The set of common positive divisors is nonempty because it contains $1$, and finite because at least one of $a,b$ is positive and bounds its divisors.

For $b>0$, integer division gives unique integers $q,r$ with
\[
a=qb+r,\qquad 0\le r<b.
\]
The integer $r$ is the remainder, also written $a\bmod b$. Here $q$ counts the whole copies of $b$ removed from $a$; $r$ is what remains. For example, $17=3\cdot5+2$, so the quotient is three and the remainder is two. The condition $0\le r<b$ identifies the permitted remainder; $17=2\cdot5+7$ is a true equality but not division with its remainder in the required range.

Euclid's algorithm is built on one idea: replace the pair $(a,b)$ by the pair $(b,r)$, where $r$ is the remainder of $a$ on division by $b$. For example, $(84,30)$ is replaced by $(30,24)$, because $84=2\cdot30+24$. The positive common divisors of $84$ and $30$ are $1,2,3,6$, and so are those of $30$ and $24$. We now show that this is no coincidence: the replacement never changes the set of common divisors.

First let $d$ divide $a$ and $b$. Then $a=ds$ and $b=dt$ for some integers $s,t$. Substitution gives
\[
 r=a-qb=ds-qdt=d(s-qt).
\]
The bracket is an integer, so $d$ divides $r$. Thus every common divisor of $a,b$ is a common divisor of $b,r$.

For the reverse direction, suppose $b=dt$ and $r=du$. Then
\[
 a=qb+r=qdt+du=d(qt+u),
\]
so $d$ divides $a$ as well. Both directions are needed: only then do we know that the two pairs have \emph{exactly} the same common divisors. The greatest common divisor is the largest member of that set of common divisors, and the two pairs have the same set, so they have the same largest member:
\[
\gcd(a,b)=\gcd(b,r).
\]
This is the invariant behind the algorithm. Repeatedly replace $(a,b)$ by $(b,a\bmod b)$ while $b>0$.

For $(a,b)=(84,30)$, the divisions are
\[
84=2\cdot30+24,\qquad30=1\cdot24+6,\qquad24=4\cdot6+0.
\]
The successive pairs are $(84,30)$, $(30,24)$, $(24,6)$, $(6,0)$. At the last pair, the greatest common divisor is $6$, because the positive divisors common to $6$ and $0$ are exactly the divisors of $6$.

Why must the process stop? Its second coordinate is a nonnegative integer and strictly decreases at every step that occurs. There cannot be an infinite strictly decreasing sequence of nonnegative integers: if the first value is $N$, each step lowers the value by at least one, so after at most $N$ steps the value would have to drop below zero. Once the second coordinate is zero, the algorithm stops and returns the first coordinate. The invariant proves the answer is right; the decreasing integer proves an answer is eventually returned.

\begin{principle}{Habit 5: separate preservation from progress}
For an iterative construction, ask two different questions. What property remains true after a step? What quantity forces progress or controls the accumulated error? A preserved quantity may prove correctness without proving termination. A decreasing real quantity may converge without ever reaching zero.
\end{principle}

\pauseq{Does a strictly decreasing nonnegative real quantity prove that a process stops?}{No. The values $1,1/2,1/4,\ldots$ decrease forever and remain positive. The Euclidean termination proof depends on the decreasing quantity being an integer, not merely a real number.}
